% Options for packages loaded elsewhere
% Options for packages loaded elsewhere
\PassOptionsToPackage{unicode}{hyperref}
\PassOptionsToPackage{hyphens}{url}
\PassOptionsToPackage{dvipsnames,svgnames,x11names}{xcolor}
%
\documentclass[
  a4paper,
  DIV=11,
  numbers=noendperiod,
  oneside]{scrreport}
\usepackage{xcolor}
\usepackage[left=10mm,marginparwidth=60mm,textwidth=130mm,marginparsep=5mm]{geometry}
\usepackage{amsmath,amssymb}
\setcounter{secnumdepth}{3}
\usepackage{iftex}
\ifPDFTeX
  \usepackage[T1]{fontenc}
  \usepackage[utf8]{inputenc}
  \usepackage{textcomp} % provide euro and other symbols
\else % if luatex or xetex
  \usepackage{unicode-math} % this also loads fontspec
  \defaultfontfeatures{Scale=MatchLowercase}
  \defaultfontfeatures[\rmfamily]{Ligatures=TeX,Scale=1}
\fi
\usepackage{lmodern}
\ifPDFTeX\else
  % xetex/luatex font selection
\fi
% Use upquote if available, for straight quotes in verbatim environments
\IfFileExists{upquote.sty}{\usepackage{upquote}}{}
\IfFileExists{microtype.sty}{% use microtype if available
  \usepackage[]{microtype}
  \UseMicrotypeSet[protrusion]{basicmath} % disable protrusion for tt fonts
}{}
\makeatletter
\@ifundefined{KOMAClassName}{% if non-KOMA class
  \IfFileExists{parskip.sty}{%
    \usepackage{parskip}
  }{% else
    \setlength{\parindent}{0pt}
    \setlength{\parskip}{6pt plus 2pt minus 1pt}}
}{% if KOMA class
  \KOMAoptions{parskip=half}}
\makeatother
% Make \paragraph and \subparagraph free-standing
\makeatletter
\ifx\paragraph\undefined\else
  \let\oldparagraph\paragraph
  \renewcommand{\paragraph}{
    \@ifstar
      \xxxParagraphStar
      \xxxParagraphNoStar
  }
  \newcommand{\xxxParagraphStar}[1]{\oldparagraph*{#1}\mbox{}}
  \newcommand{\xxxParagraphNoStar}[1]{\oldparagraph{#1}\mbox{}}
\fi
\ifx\subparagraph\undefined\else
  \let\oldsubparagraph\subparagraph
  \renewcommand{\subparagraph}{
    \@ifstar
      \xxxSubParagraphStar
      \xxxSubParagraphNoStar
  }
  \newcommand{\xxxSubParagraphStar}[1]{\oldsubparagraph*{#1}\mbox{}}
  \newcommand{\xxxSubParagraphNoStar}[1]{\oldsubparagraph{#1}\mbox{}}
\fi
\makeatother

\usepackage{color}
\usepackage{fancyvrb}
\newcommand{\VerbBar}{|}
\newcommand{\VERB}{\Verb[commandchars=\\\{\}]}
\DefineVerbatimEnvironment{Highlighting}{Verbatim}{commandchars=\\\{\}}
% Add ',fontsize=\small' for more characters per line
\usepackage{framed}
\definecolor{shadecolor}{RGB}{241,243,245}
\newenvironment{Shaded}{\begin{snugshade}}{\end{snugshade}}
\newcommand{\AlertTok}[1]{\textcolor[rgb]{0.68,0.00,0.00}{#1}}
\newcommand{\AnnotationTok}[1]{\textcolor[rgb]{0.37,0.37,0.37}{#1}}
\newcommand{\AttributeTok}[1]{\textcolor[rgb]{0.40,0.45,0.13}{#1}}
\newcommand{\BaseNTok}[1]{\textcolor[rgb]{0.68,0.00,0.00}{#1}}
\newcommand{\BuiltInTok}[1]{\textcolor[rgb]{0.00,0.23,0.31}{#1}}
\newcommand{\CharTok}[1]{\textcolor[rgb]{0.13,0.47,0.30}{#1}}
\newcommand{\CommentTok}[1]{\textcolor[rgb]{0.37,0.37,0.37}{#1}}
\newcommand{\CommentVarTok}[1]{\textcolor[rgb]{0.37,0.37,0.37}{\textit{#1}}}
\newcommand{\ConstantTok}[1]{\textcolor[rgb]{0.56,0.35,0.01}{#1}}
\newcommand{\ControlFlowTok}[1]{\textcolor[rgb]{0.00,0.23,0.31}{\textbf{#1}}}
\newcommand{\DataTypeTok}[1]{\textcolor[rgb]{0.68,0.00,0.00}{#1}}
\newcommand{\DecValTok}[1]{\textcolor[rgb]{0.68,0.00,0.00}{#1}}
\newcommand{\DocumentationTok}[1]{\textcolor[rgb]{0.37,0.37,0.37}{\textit{#1}}}
\newcommand{\ErrorTok}[1]{\textcolor[rgb]{0.68,0.00,0.00}{#1}}
\newcommand{\ExtensionTok}[1]{\textcolor[rgb]{0.00,0.23,0.31}{#1}}
\newcommand{\FloatTok}[1]{\textcolor[rgb]{0.68,0.00,0.00}{#1}}
\newcommand{\FunctionTok}[1]{\textcolor[rgb]{0.28,0.35,0.67}{#1}}
\newcommand{\ImportTok}[1]{\textcolor[rgb]{0.00,0.46,0.62}{#1}}
\newcommand{\InformationTok}[1]{\textcolor[rgb]{0.37,0.37,0.37}{#1}}
\newcommand{\KeywordTok}[1]{\textcolor[rgb]{0.00,0.23,0.31}{\textbf{#1}}}
\newcommand{\NormalTok}[1]{\textcolor[rgb]{0.00,0.23,0.31}{#1}}
\newcommand{\OperatorTok}[1]{\textcolor[rgb]{0.37,0.37,0.37}{#1}}
\newcommand{\OtherTok}[1]{\textcolor[rgb]{0.00,0.23,0.31}{#1}}
\newcommand{\PreprocessorTok}[1]{\textcolor[rgb]{0.68,0.00,0.00}{#1}}
\newcommand{\RegionMarkerTok}[1]{\textcolor[rgb]{0.00,0.23,0.31}{#1}}
\newcommand{\SpecialCharTok}[1]{\textcolor[rgb]{0.37,0.37,0.37}{#1}}
\newcommand{\SpecialStringTok}[1]{\textcolor[rgb]{0.13,0.47,0.30}{#1}}
\newcommand{\StringTok}[1]{\textcolor[rgb]{0.13,0.47,0.30}{#1}}
\newcommand{\VariableTok}[1]{\textcolor[rgb]{0.07,0.07,0.07}{#1}}
\newcommand{\VerbatimStringTok}[1]{\textcolor[rgb]{0.13,0.47,0.30}{#1}}
\newcommand{\WarningTok}[1]{\textcolor[rgb]{0.37,0.37,0.37}{\textit{#1}}}

\usepackage{longtable,booktabs,array}
\usepackage{calc} % for calculating minipage widths
% Correct order of tables after \paragraph or \subparagraph
\usepackage{etoolbox}
\makeatletter
\patchcmd\longtable{\par}{\if@noskipsec\mbox{}\fi\par}{}{}
\makeatother
% Allow footnotes in longtable head/foot
\IfFileExists{footnotehyper.sty}{\usepackage{footnotehyper}}{\usepackage{footnote}}
\makesavenoteenv{longtable}
\usepackage{graphicx}
\makeatletter
\newsavebox\pandoc@box
\newcommand*\pandocbounded[1]{% scales image to fit in text height/width
  \sbox\pandoc@box{#1}%
  \Gscale@div\@tempa{\textheight}{\dimexpr\ht\pandoc@box+\dp\pandoc@box\relax}%
  \Gscale@div\@tempb{\linewidth}{\wd\pandoc@box}%
  \ifdim\@tempb\p@<\@tempa\p@\let\@tempa\@tempb\fi% select the smaller of both
  \ifdim\@tempa\p@<\p@\scalebox{\@tempa}{\usebox\pandoc@box}%
  \else\usebox{\pandoc@box}%
  \fi%
}
% Set default figure placement to htbp
\def\fps@figure{htbp}
\makeatother





\setlength{\emergencystretch}{3em} % prevent overfull lines

\providecommand{\tightlist}{%
  \setlength{\itemsep}{0pt}\setlength{\parskip}{0pt}}



 


% This is used for pdf compilation from qmd.
% Any special LaTeX commands should be updated in parallel with mymacros.sty which is used for html compilation.
% All rows *might* need an HTML equivalent. Those at the bottom of this file are identical.
% This is how the two files are included in the _quarto.yaml header of the qmd file
%format:
%  html:
%    defaults: [resources/latex/global-maths-definitions.yaml]
%  pdf:
%   include-in-header: ["include-in-header.tex"]

% the 'defaults' above secretly points to mymacros.sty where they really live.

\usepackage{gensymb}
\usepackage{amsmath}
\usepackage{xcolor}

\usepackage[skins,breakable]{tcolorbox}

% Code to reduce width of code blocks in PDF to encourage line wrapping
\usepackage{fvextra}
      \DefineVerbatimEnvironment{Highlighting}{Verbatim}{
        breaklines,
        breakanywhere,
        commandchars=\\\{\},
        fontsize=\small
      }
% Also line wrap in verbatim blocks
\RecustomVerbatimEnvironment{verbatim}{Verbatim}{breaklines,fontsize=\footnotesize}

\usepackage{environ}

% Below here is an old way to defining new custom coloured boxes.
% We use nice extensions now.

% \definecolor{goals-col}{RGB}{201, 232, 222}
% 
% % For defining new callout boxes
% % HTML version is in CSS file
% % Need to turn on goals and goals-header as environments using the latex-environment extension for quarto
% % This is designed to work for the following fenced div syntax:
% % :::{.goals}
% % ::::{.goals-header}
% % TITLE
% % ::::
% % Actual box contents
% % :::
% \newtcolorbox{goals}[1][]{notitle, top=0pt, boxsep=0pt, coltitle=black!90!white, colback=white, colframe=goals-col,#1}
% \NewEnviron{goals-header}{\tcbsubtitle[top=1mm, bottom=1mm]{\BODY}}
% % It works by creating a box with no title. Then no spacing. Then a subtitle defined later.

% Increasing spacing between rows of tables
\def\arraystretch{2}

% Syntax for highlighting my KEY TERMS
% HTML version is included in css for .term
\newcommand{\term}[1]{\textbf{#1}}

% Added to do derivatives, custom commands for HTML
% read by MathJax are inserted via diff.macro
% must keep naming in sync with esdiff package
\usepackage{esdiff}

%%%%%%%%%%%%%%%%%%%%%%%%%%%%%%%%%%%%%%%%%%%%%%%%%%%%%%%%%%%%
%%% BELOW HERE ARE COMMANDS SHARED WITH THE HTML OUTPUT %%%%
%%%%%%%%%%%%%%%%%%%%%%%%%%%%%%%%%%%%%%%%%%%%%%%%%%%%%%%%%%%%
%%%% KEEP THESE UPDATED ALSO IN THE MYMACROS.STY FILE   %%%%
%%%%%%%%%%%%%%%%%%%%%%%%%%%%%%%%%%%%%%%%%%%%%%%%%%%%%%%%%%%%

% Adding some operator names missing by default
\DeclareMathOperator{\cosec}{cosec}
\DeclareMathOperator{\powop}{pow} %\newcommand{\pow}{\operatorname{pow}}
\definecolor{frog}{rgb}{1,0,0}
\newcommand{\dx}{\,\text{d}x}
\newcommand{\du}{\,\text{d}u}
\newcommand{\dy}{\,\text{d}y}
\newcommand{\dt}{\,\text{d}t}

\KOMAoption{captions}{tableheading}
\makeatletter
\@ifpackageloaded{float}{}{\usepackage{float}}
\floatstyle{plain}
\@ifundefined{c@chapter}{\newfloat{video}{h}{lovid}}{\newfloat{video}{h}{lovid}[chapter]}
\floatname{video}{Video}
\newcommand*\listofvideos{\listof{video}{List of Videos}}
\makeatother
\makeatletter
\@ifpackageloaded{tcolorbox}{}{\usepackage[skins,breakable]{tcolorbox}}
\@ifpackageloaded{fontawesome5}{}{\usepackage{fontawesome5}}
\definecolor{quarto-callout-color}{HTML}{909090}
\definecolor{quarto-callout-note-color}{HTML}{0758E5}
\definecolor{quarto-callout-important-color}{HTML}{CC1914}
\definecolor{quarto-callout-warning-color}{HTML}{EB9113}
\definecolor{quarto-callout-tip-color}{HTML}{00A047}
\definecolor{quarto-callout-caution-color}{HTML}{FC5300}
\definecolor{quarto-callout-color-frame}{HTML}{acacac}
\definecolor{quarto-callout-note-color-frame}{HTML}{4582ec}
\definecolor{quarto-callout-important-color-frame}{HTML}{d9534f}
\definecolor{quarto-callout-warning-color-frame}{HTML}{f0ad4e}
\definecolor{quarto-callout-tip-color-frame}{HTML}{02b875}
\definecolor{quarto-callout-caution-color-frame}{HTML}{fd7e14}
\makeatother
\makeatletter
\@ifpackageloaded{bookmark}{}{\usepackage{bookmark}}
\makeatother
\makeatletter
\@ifpackageloaded{caption}{}{\usepackage{caption}}
\AtBeginDocument{%
\ifdefined\contentsname
  \renewcommand*\contentsname{Table of contents}
\else
  \newcommand\contentsname{Table of contents}
\fi
\ifdefined\listfigurename
  \renewcommand*\listfigurename{List of Figures}
\else
  \newcommand\listfigurename{List of Figures}
\fi
\ifdefined\listtablename
  \renewcommand*\listtablename{List of Tables}
\else
  \newcommand\listtablename{List of Tables}
\fi
\ifdefined\figurename
  \renewcommand*\figurename{Figure}
\else
  \newcommand\figurename{Figure}
\fi
\ifdefined\tablename
  \renewcommand*\tablename{Table}
\else
  \newcommand\tablename{Table}
\fi
}
\@ifpackageloaded{float}{}{\usepackage{float}}
\floatstyle{ruled}
\@ifundefined{c@chapter}{\newfloat{codelisting}{h}{lop}}{\newfloat{codelisting}{h}{lop}[chapter]}
\floatname{codelisting}{Listing}
\newcommand*\listoflistings{\listof{codelisting}{List of Listings}}
\makeatother
\makeatletter
\makeatother
\makeatletter
\@ifpackageloaded{caption}{}{\usepackage{caption}}
\@ifpackageloaded{subcaption}{}{\usepackage{subcaption}}
\makeatother
\makeatletter
\@ifpackageloaded{sidenotes}{}{\usepackage{sidenotes}}
\@ifpackageloaded{marginnote}{}{\usepackage{marginnote}}
\makeatother
\makeatletter
\@ifpackageloaded{tcolorbox}{}{\usepackage[many]{tcolorbox}}
\makeatother
%%%% ---foldboxy preamble ----- %%%%%

\definecolor{fbx-default-color1}{HTML}{c7c7d0}
\definecolor{fbx-default-color2}{HTML}{a3a3aa}

\definecolor{fbox-color1}{HTML}{c7c7d0}
\definecolor{fbox-color2}{HTML}{a3a3aa}

% arguments: #1 typelabelnummer: #2 titel: #3
\newenvironment{fbx}[3]{\begin{tcolorbox}[enhanced, breakable,%
attach boxed title to top*={xshift=1.4pt},
boxed title style={boxrule=0.0mm, fuzzy shadow={1pt}{-1pt}{0mm}{0.1mm}{gray}, arc=.3em, rounded corners=east, sharp corners=west}, colframe=#1-color2, colbacktitle=#1-color1, colback = white, coltitle=black,  titlerule=0mm, toprule=0pt, bottomrule=.7pt, leftrule=.3em, rightrule=0pt, outer arc=.3em,  arc=0pt,	 sharp corners = east, left=.5em, bottomtitle=1mm, toptitle=1mm,title=\textbf{#2}\hspace{0.5em}{#3}]}
{\end{tcolorbox}}

% boxed environment with right border
\newenvironment{fbxSimple}[3]{\begin{tcolorbox}[enhanced, breakable,%
attach boxed title to top*={xshift=1.4pt},
boxed title style={boxrule=0.0mm, fuzzy shadow={1pt}{-1pt}{0mm}{0.1mm}{gray}, arc=.3em, rounded corners=east, sharp corners=west}, colframe=#1-color2, colbacktitle=#1-color1, colback = white, coltitle=black,  titlerule=0mm, toprule=0pt, bottomrule=.7pt, leftrule=.3em, rightrule=.7pt, outer arc=.3em,  	left=.5em, right=.5em, bottomtitle=1mm, toptitle=1mm,title=\textbf{#2}\hspace{0.5em}{#3}]}
{\end{tcolorbox}}

%%%% --- end foldboxy preamble ----- %%%%%
%%==== colors from yaml ===%
\definecolor{Answer-color1}{HTML}{d18888}
\definecolor{Answer-color2}{HTML}{660000}
\definecolor{Lemma-color1}{HTML}{b2dac4}
\definecolor{Lemma-color2}{HTML}{004721}
\definecolor{Corollary-color1}{HTML}{7fbad7}
\definecolor{Corollary-color2}{HTML}{005398}
\definecolor{Video-color1}{HTML}{d18888}
\definecolor{Video-color2}{HTML}{660000}
\definecolor{Theorem-color1}{HTML}{948bde}
\definecolor{Theorem-color2}{HTML}{584eab}
\definecolor{Example-color1}{HTML}{b2dac4}
\definecolor{Example-color2}{HTML}{004721}
\definecolor{Definition-color1}{HTML}{ffdc36}
\definecolor{Definition-color2}{HTML}{ffb948}
\definecolor{Task-color1}{HTML}{948bde}
\definecolor{Task-color2}{HTML}{584eab}
\definecolor{Weblink-color1}{HTML}{d18888}
\definecolor{Weblink-color2}{HTML}{660000}
%=============%
\newcounter{quartocallouttipno}
\newcommand{\quartocallouttip}[1]{\refstepcounter{quartocallouttipno}\label{#1}}
\newcounter{quartocalloutcauno}
\newcommand{\quartocalloutcau}[1]{\refstepcounter{quartocalloutcauno}\label{#1}}
\usepackage{bookmark}
\IfFileExists{xurl.sty}{\usepackage{xurl}}{} % add URL line breaks if available
\urlstyle{same}
\hypersetup{
  colorlinks=true,
  linkcolor={blue},
  filecolor={Maroon},
  citecolor={Blue},
  urlcolor={Blue},
  pdfcreator={LaTeX via pandoc}}


\author{}
\date{}
\begin{document}

\renewcommand*\contentsname{Table of contents}
{
\hypersetup{linkcolor=}
\setcounter{tocdepth}{2}
\tableofcontents
}

\bookmarksetup{startatroot}

\chapter*{Welcome to this book of
notes}\label{welcome-to-this-book-of-notes}
\addcontentsline{toc}{chapter}{Welcome to this book of notes}

\markboth{Welcome to this book of notes}{Welcome to this book of notes}

\bookmarksetup{startatroot}

\chapter{Chapter 1 - Headings}\label{sec-headings}

Each row that begins with a \# is a chapter heading. So if each chapter
is a single file it's best to only use a single \# line at the top of
each document.

\section{Models for binary/binomial response}\label{sec-binary-response}

This chapter contains sections and subsections. This is the intro. Since
the chapter title begins with a \texttt{\#}, sections begin with
\texttt{\#\#}, and subsections should begin with \texttt{\#\#\#}.

Typical default numbering will look like this:

\begin{verbatim}
Chapter 1
Section 1.1
Subsection 1.1.1
Subsection 1.1.2
Section 1.2
Subsection 1.2.1
Subsection 1.2.2
Subsection 1.2.3
Section 1.3
\end{verbatim}

The \texttt{crossref} option in your global \texttt{.yml} file does
offer the chance to not number any chapters and just have sections in
each chapter, if you really don't want chapters.

\subsection{First subsection}\label{first-subsection}

Text goes here.

\subsection{Second subsection}\label{second-subsection}

Text goes here. And we can reference whole sections if we wish, like
Section~\ref{sec-binary-response}. Indeed if you look in
Chapter~\ref{sec-links}.

Let

\[D_3(y) \equiv \diff[3]{y}{x}\]

That was just an equation inside \texttt{\$\$} and \texttt{\$\$} tags.
For labelled equations you need to do the slightly unusual LaTeX
behaviour of using \texttt{\$\$} and then
\texttt{\textbackslash{}\textbackslash{}begin\{aligned\}} or similar
inside. So unlike in LaTeX you always need \texttt{\$\$} to start a math
environment (because the \texttt{\$\$} tells the browser to enable math
mode). Yes! No.~Maybe. Not.

\begin{equation}\phantomsection\label{eq-AROC}{
\begin{aligned}
\int_0^{\infty} x^3 \dx
\end{aligned}
}\end{equation}

This week we will learn how to model outcomes of interest that take one
of two categorical values (e.g.~yes/no, success/failure, alive/dead).
The independent responses \(Y_i\) can either be

\begin{itemize}
\item
  \emph{binary} (ungrouped), taking the value 1 (say success, with
  probability \(p_i\)) or 0 (failure, with probability \(1-p_i\)) or
\item
  \emph{binomial} (grouped), where \(Y_i\) is the number of successes in
  a given number of trials \(n_i\), with the probability of success
  being \(p_i\) and the probability of failure being \(1-p_i\).
\end{itemize}

In both cases the distribution of the \(Y_i\) is binomial, but in the
first case it is Bin\((1,p_i)\) and in the second case it is
Bin\((n_i,p_i)\).

\section{Binomial response}\label{binomial-response}

We will begin with models for binomial responses and we will look at
exploratory plots of the data, different choices of link function and
hypothesis tests about terms in the model. Finally we will examine
measures of goodness of fit of the model. Let's begin with an example.

\begin{tcolorbox}[enhanced jigsaw, coltitle=black, leftrule=.75mm, bottomtitle=1mm, colframe=quarto-callout-tip-color-frame, opacitybacktitle=0.6, colback=white, toprule=.15mm, colbacktitle=quarto-callout-tip-color!10!white, toptitle=1mm, titlerule=0mm, title=\textcolor{quarto-callout-tip-color}{\faLightbulb}\hspace{0.5em}{Beetle mortality}, arc=.35mm, opacityback=0, rightrule=.15mm, bottomrule=.15mm, breakable, left=2mm]

This is a ``Tip Callout'' box called via

\begin{verbatim}
:::{.callout-tip}
### Banner title like Bettle Mortality

Content here. This gives coloured boxes in HTML and LaTeX.
:::
\end{verbatim}

Now for a default code \texttt{panelset} as Quarto calls it. For
allowing dual code presentation. See Chapter~\ref{sec-code-demo} for my
restyled version.

\subsection{R}

\begin{Shaded}
\begin{Highlighting}[]
\NormalTok{beetles }\OtherTok{\textless{}{-}} \FunctionTok{read.csv}\NormalTok{(}\FunctionTok{url}\NormalTok{(}\StringTok{"http://www.stats.gla.ac.uk/\textasciitilde{}tereza/rp/beetles.csv"}\NormalTok{))}
\NormalTok{beetles}
\end{Highlighting}
\end{Shaded}

\subsection{Python}

\begin{Shaded}
\begin{Highlighting}[]
\ImportTok{import}\NormalTok{ pandas }\ImportTok{as}\NormalTok{ pd}
\CommentTok{\# beetles = pd.read\_csv("beetles.csv")}
\CommentTok{\# beetles}
\end{Highlighting}
\end{Shaded}

\end{tcolorbox}

\begin{tcolorbox}[enhanced jigsaw, coltitle=black, leftrule=.75mm, bottomtitle=1mm, colframe=quarto-callout-warning-color-frame, opacitybacktitle=0.6, colback=white, toprule=.15mm, colbacktitle=quarto-callout-warning-color!10!white, toptitle=1mm, titlerule=0mm, title=\textcolor{quarto-callout-warning-color}{\faExclamationTriangle}\hspace{0.5em}{Warning \ref*{tip-special}: Example}, arc=.35mm, opacityback=0, rightrule=.15mm, bottomrule=.15mm, breakable, left=2mm]

\quartocallouttip{tip-special} 

Some text

\end{tcolorbox}

This box above can be linked to in the HTML via Tip~\ref{tip-special}.

\begin{tcolorbox}[enhanced jigsaw, coltitle=black, leftrule=.75mm, bottomtitle=1mm, colframe=quarto-callout-warning-color-frame, opacitybacktitle=0.6, colback=white, toprule=.15mm, colbacktitle=quarto-callout-warning-color!10!white, toptitle=1mm, titlerule=0mm, title=\textcolor{quarto-callout-warning-color}{\faExclamationTriangle}\hspace{0.5em}{Example}, arc=.35mm, opacityback=0, rightrule=.15mm, bottomrule=.15mm, breakable, left=2mm]

Some text

\end{tcolorbox}

\begin{tcolorbox}[enhanced jigsaw, coltitle=black, leftrule=.75mm, bottomtitle=1mm, colframe=quarto-callout-warning-color-frame, opacitybacktitle=0.6, colback=white, toprule=.15mm, colbacktitle=quarto-callout-warning-color!10!white, toptitle=1mm, titlerule=0mm, title=\textcolor{quarto-callout-warning-color}{\faExclamationTriangle}\hspace{0.5em}{Example}, arc=.35mm, opacityback=0, rightrule=.15mm, bottomrule=.15mm, breakable, left=2mm]

Some text

\end{tcolorbox}

\begin{tcolorbox}[enhanced jigsaw, coltitle=black, leftrule=.75mm, bottomtitle=1mm, colframe=quarto-callout-warning-color-frame, opacitybacktitle=0.6, colback=white, toprule=.15mm, colbacktitle=quarto-callout-warning-color!10!white, toptitle=1mm, titlerule=0mm, title=\textcolor{quarto-callout-warning-color}{\faExclamationTriangle}\hspace{0.5em}{Warning \ref*{tip-special2}: Example}, arc=.35mm, opacityback=0, rightrule=.15mm, bottomrule=.15mm, breakable, left=2mm]

\quartocallouttip{tip-special2} 

Some text

\end{tcolorbox}

Only the first and last of these had a label. They were
Tip~\ref{tip-special} and Tip~\ref{tip-special2}. Notice, however, that
you will need to define your own custom blocks. Since a \texttt{warning}
callout, which is used for an ``Example'' will suddenly be called
``Warning 1.1'' since the default naming of a labelled (and thus
numbered) block would need to be overridden.

Also notice the massive confusion caused by me using the
\texttt{\#tip-name} style label on a \texttt{warning} callout! Make sure
to use the right prefix in the reference else it'll do this.

\begin{tcolorbox}[enhanced jigsaw, coltitle=black, leftrule=.75mm, bottomtitle=1mm, colframe=quarto-callout-caution-color-frame, opacitybacktitle=0.6, colback=white, toprule=.15mm, colbacktitle=quarto-callout-caution-color!10!white, toptitle=1mm, titlerule=0mm, title=\textcolor{quarto-callout-caution-color}{\faFire}\hspace{0.5em}{Caution \ref*{cau-correct}: this is a caution callout}, arc=.35mm, opacityback=0, rightrule=.15mm, bottomrule=.15mm, breakable, left=2mm]

\quartocalloutcau{cau-correct} 

The reference name begins with \texttt{\#cau-}.

\end{tcolorbox}

Then when I reference it like this: Caution~\ref{cau-correct}, all is
well.

\bookmarksetup{startatroot}

\chapter{Week 2 -- Code demonstrations}\label{sec-code-demo}

Some equations, again. Notice this chapter begins with:

\begin{verbatim}
---
filters:
    - acc-tools
---
\end{verbatim}

This loads the bonus accessibility features extension I have made.

\textbf{LET'S MAKE THIS CHAPTER ALL ABOUT SHOWING OFF THE ACCESSIBILITY
STUFF}

\[\diff[3]{y}{x}\]

\[\int_0^{\infty} x^2 \dx\]

This equation is the same as Equation~\ref{eq-AROC} from Chapter 1.

\section{Features}\label{features}

\begin{itemize}
\tightlist
\item
  You'll first note in the bottom left is a global toggle button set.
  You can change all code blocks to R or to Python, by clicking; or by
  pressing \texttt{R} or \texttt{P} on the keyboard.
\end{itemize}

\section{Binomial response}\label{binomial-response-1}

We will begin with models for binomial responses and we will look at
exploratory plots of the data, different choices of link function and
hypothesis tests about terms in the model. Finally we will examine
measures of goodness of fit of the model. Let's begin with an example.

\begin{tcolorbox}[enhanced jigsaw, coltitle=black, leftrule=.75mm, bottomtitle=1mm, colframe=quarto-callout-tip-color-frame, opacitybacktitle=0.6, colback=white, toprule=.15mm, colbacktitle=quarto-callout-tip-color!10!white, toptitle=1mm, titlerule=0mm, title=\textcolor{quarto-callout-tip-color}{\faLightbulb}\hspace{0.5em}{Beetle mortality}, arc=.35mm, opacityback=0, rightrule=.15mm, bottomrule=.15mm, breakable, left=2mm]

Jump to the \hyperref[beetle-code]{Beetle code}.

Data for this example consists of the number of beetles dead
(\emph{killed}) after five hours exposure to gaseous carbon disulphide
at various concentrations (\emph{dose}). The goal for this analysis is
to model the probability of a beetle dying as a function of the carbon
disulphide dose.

\subsection{R}

\begin{Shaded}
\begin{Highlighting}[]
\NormalTok{beetles }\OtherTok{\textless{}{-}} \FunctionTok{read.csv}\NormalTok{(}\FunctionTok{url}\NormalTok{(}\StringTok{"http://www.stats.gla.ac.uk/\textasciitilde{}tereza/rp/beetles.csv"}\NormalTok{))}
\NormalTok{beetles}
\end{Highlighting}
\end{Shaded}

\subsection{Python}

\begin{Shaded}
\begin{Highlighting}[]
\ImportTok{import}\NormalTok{ pandas }\ImportTok{as}\NormalTok{ pd}
\CommentTok{\# beetles = pd.read\_csv("resources/data/beetles.csv")}
\CommentTok{\# beetles}
\end{Highlighting}
\end{Shaded}

\end{tcolorbox}

\begin{tcolorbox}[enhanced jigsaw, coltitle=black, leftrule=.75mm, bottomtitle=1mm, colframe=quarto-callout-tip-color-frame, opacitybacktitle=0.6, colback=white, toprule=.15mm, colbacktitle=quarto-callout-tip-color!10!white, toptitle=1mm, titlerule=0mm, title=\textcolor{quarto-callout-tip-color}{\faLightbulb}\hspace{0.5em}{Conventional tip callout}, arc=.35mm, opacityback=0, rightrule=.15mm, bottomrule=.15mm, breakable, left=2mm]

Some content

\end{tcolorbox}

\begin{tcolorbox}[enhanced jigsaw, coltitle=black, leftrule=.75mm, bottomtitle=1mm, colframe=quarto-callout-warning-color-frame, opacitybacktitle=0.6, colback=white, toprule=.15mm, colbacktitle=quarto-callout-warning-color!10!white, toptitle=1mm, titlerule=0mm, title=\textcolor{quarto-callout-warning-color}{\faExclamationTriangle}\hspace{0.5em}{Conventional warning callout}, arc=.35mm, opacityback=0, rightrule=.15mm, bottomrule=.15mm, breakable, left=2mm]

Some content

\end{tcolorbox}

\begin{tcolorbox}[enhanced jigsaw, coltitle=black, leftrule=.75mm, bottomtitle=1mm, colframe=quarto-callout-caution-color-frame, opacitybacktitle=0.6, colback=white, toprule=.15mm, colbacktitle=quarto-callout-caution-color!10!white, toptitle=1mm, titlerule=0mm, title=\textcolor{quarto-callout-caution-color}{\faFire}\hspace{0.5em}{Conventional caution callout}, arc=.35mm, opacityback=0, rightrule=.15mm, bottomrule=.15mm, breakable, left=2mm]

Some content

\end{tcolorbox}

\begin{tcolorbox}[enhanced jigsaw, coltitle=black, leftrule=.75mm, bottomtitle=1mm, colframe=quarto-callout-note-color-frame, opacitybacktitle=0.6, colback=white, toprule=.15mm, colbacktitle=quarto-callout-note-color!10!white, toptitle=1mm, titlerule=0mm, title=\textcolor{quarto-callout-note-color}{\faInfo}\hspace{0.5em}{Conventional note callout}, arc=.35mm, opacityback=0, rightrule=.15mm, bottomrule=.15mm, breakable, left=2mm]

Some content

\end{tcolorbox}

\bookmarksetup{startatroot}

\chapter{Links chapter}\label{sec-links}

This chapter can even link back to other whole chapters
Section~\ref{sec-binary-response}. Or equations in other chapters
Equation~\ref{eq-AROC}.

\section{Section One}\label{section-one}

One way to include videos is to use the Quarto syntax used to insert any
figures, where you write

\begin{verbatim}
::: {#vid-vid3}
{{﹤video https://youtu.be/d2nG7Y17sQw 
  title = "Binomial response models applied to beetle mortality data (9m59s)",
  <other customizations can be put here>
>}}

This is the caption.
:::
\end{verbatim}

For more details:
\href{https://quarto.org/docs/authoring/videos.html}{Quarto video
manual}

Here it is:

\begin{video}

\centering{

\url{https://youtu.be/d2nG7Y17sQw}

}

\caption{\label{vid-vid3}This is the caption.}

\end{video}%

Referencing is simple, like this Video~\ref{vid-vid3}. Such video blocks
must contain actual text (the caption is drawn from the final line).

\section{Equations}\label{sec-equations}

There is good Quarto documentation on how to crossref things
\href{https://quarto.org/docs/authoring/cross-references.html}{Quarto
Cross References}.

The key thing to note is that if you plan to references a standard
object like a heading/section the internal label name is restricted to
begin with a particular three-letter prefix (except equations which get
just \texttt{eq}). Plus they recommend against using underscores in
labels.

New prefixes can be created, as this template does with to create new
figures for videos called \texttt{vid}. Theorems, Lemmas etc.. this
template recommends the numbered-custom-boxes demonstrated in another
chapter. If you follow this recommendation then you must avoid using
names like \texttt{\#thm-fermat} because Quarto will be confused about
which method you wish. So in the examples presented I have just started
all references with \texttt{\#r} to avoid overlap,
e.g.~\texttt{\#rthm-fermat}.

Here's an equation: \begin{equation}\phantomsection\label{eq-squares}{
\begin{aligned}
x^2 + y^2 = z^2
\end{aligned}
}\end{equation} This equation was created using Quarto labelling syntax,
i.e.~like this

\begin{verbatim}
$$
\begin{aligned}
x^2 + y^2 = z^2
\end{aligned}
$$ {#eq-squares}
\end{verbatim}

Then you type \texttt{@eq-squares} to get Equation~\ref{eq-squares}.

This section was defined like this

\begin{verbatim}
## Equations {#sec-equations}
\end{verbatim}

as such we can now type \texttt{@sec-equations} and it will render as
Section~\ref{sec-equations}.

Quarto even provides a clever `hover to see a preview function'. If you
don't want it there are ways to disable it (I found it messed with
tabbing through the page on one of my machines).

The main drawback to using the Quarto numbering system is for equation
blocks (equations on many consecutive lines), where the way LaTeX offers
individually labelling of lines Quarto doesn't. There is a more complex
workaround to delegate all labelling of equations to MathJax (the
javascript equation renderer but I haven't gone with that).

\bookmarksetup{startatroot}

\chapter{Chapter 4 - Nicer boxes}\label{chapter-4---nicer-boxes}

This chapter shows off some nice coloured boxes. They come from the
\texttt{custom-numbered-blocks} extension. First you need to install the
extension \texttt{quarto\ add\ ute/custom-numbered-blocks}.

This code has then been added to the \texttt{\_quarto.yml} or
\texttt{\_metadata.yml} file (or at the top of this file) to load the
coloured boxes.

\begin{verbatim}
filters:
  - custom-numbered-blocks # This loads the extension

# This is ute-hahn's nice coloured boxes: setup
custom-numbered-blocks:
  groups:
    thmlike:
      colors: [948bde, 584eab]
      boxstyle: foldbox.simple
    deflike:
      colors: [ffdc36, ffb948]
      boxstyle: foldbox.simple
    videolike:
      colors: [beb9eb, 877eda]
      boxstyle: foldbox.simple
      collapse: open
  classes:
    Theorem:
      group: thmlike
    Lemma:
      group: thmlike
    Corollary:
      group: thmlike
    Definition:
      group: deflike
    Video:
      group: videolike
      
crossref:
  chapters: true
\end{verbatim}

\section{Custom Numbered Blocks}\label{custom-numbered-blocks}

We've defined several custom-numbered block types using the
\texttt{custom-numbered-blocks} option. The key point is that these can
be numbered and follow grouped numbering conventions.

Each block belongs to a \textbf{group} that determines its shared
numbering.

All of these blocks are \emph{foldable}, meaning they can be opened and
closed. The \texttt{collapse} setting determines the starting state.

\begin{longtable}[]{@{}
  >{\raggedright\arraybackslash}p{(\linewidth - 8\tabcolsep) * \real{0.1429}}
  >{\raggedright\arraybackslash}p{(\linewidth - 8\tabcolsep) * \real{0.1039}}
  >{\raggedright\arraybackslash}p{(\linewidth - 8\tabcolsep) * \real{0.2338}}
  >{\raggedright\arraybackslash}p{(\linewidth - 8\tabcolsep) * \real{0.3117}}
  >{\raggedright\arraybackslash}p{(\linewidth - 8\tabcolsep) * \real{0.2078}}@{}}
\caption{This is the table's caption}\label{tbl-boxcols}\tabularnewline
\toprule\noalign{}
\begin{minipage}[b]{\linewidth}\raggedright
Box Name
\end{minipage} & \begin{minipage}[b]{\linewidth}\raggedright
Group
\end{minipage} & \begin{minipage}[b]{\linewidth}\raggedright
Title Background Color
\end{minipage} & \begin{minipage}[b]{\linewidth}\raggedright
Left Border Color
\end{minipage} & \begin{minipage}[b]{\linewidth}\raggedright
Default State
\end{minipage} \\
\midrule\noalign{}
\endfirsthead
\toprule\noalign{}
\begin{minipage}[b]{\linewidth}\raggedright
Box Name
\end{minipage} & \begin{minipage}[b]{\linewidth}\raggedright
Group
\end{minipage} & \begin{minipage}[b]{\linewidth}\raggedright
Title Background Color
\end{minipage} & \begin{minipage}[b]{\linewidth}\raggedright
Left Border Color
\end{minipage} & \begin{minipage}[b]{\linewidth}\raggedright
Default State
\end{minipage} \\
\midrule\noalign{}
\endhead
\bottomrule\noalign{}
\endlastfoot
Theorem & \texttt{thmlike} & {}\texttt{\#948bde} & {}\texttt{\#584eab} &
Open / Closeable \\
Lemma & \texttt{thmlike} & {}\texttt{\#b2dac4} & {}\texttt{\#004721} &
Open / Closeable \\
Corollary & \texttt{thmlike} & {}\texttt{\#7fbad7} & {}\texttt{\#005398}
& Open / Closeable \\
Definition & \texttt{deflike} & {}\texttt{\#ffdc36} &
{}\texttt{\#ffb948} & Open / Closeable \\
Video & \texttt{videolike} & {}\texttt{\#d18888} & {}\texttt{\#660000} &
Open / Closeable \\
\end{longtable}

These boxes also have dark-mode versions, see below. Tables like this
can have captions, and be referenced Table~\ref{tbl-boxcols}. They can
also have styling added.

\begin{quote}
Note: The two colors defined for each group represent the \textbf{left
border} and the \textbf{title background}.
\end{quote}

Here's an alternative to treating videos are float/figures, put them in
their own box and use the title of the box to describe the video.

\phantomsection\label{vidFirst}
\begin{fbxSimple}{Video}{Video 1: }{Binomial response models applied to beetle mortality data (9m59)}
\phantomsection\label{vidFirst}
There is some text here before the video

\url{https://youtu.be/d2nG7Y17sQw}

\end{fbxSimple}

\subsection{Beetle mortality data}\label{beetle-mortality-data}

Here's some famous beetle data. The R code is folded (and thus
hideable), the Python code isn't folded so cannot be folded away.

\subsection{R}

\begin{Shaded}
\begin{Highlighting}[]
\NormalTok{beetles }\OtherTok{\textless{}{-}} \FunctionTok{read.csv}\NormalTok{(}\FunctionTok{url}\NormalTok{(}\StringTok{"http://www.stats.gla.ac.uk/\textasciitilde{}tereza/rp/beetles.csv"}\NormalTok{))}
\NormalTok{beetles}
\end{Highlighting}
\end{Shaded}

\subsection{Python}

\begin{Shaded}
\begin{Highlighting}[]
\ImportTok{import}\NormalTok{ pandas }\ImportTok{as}\NormalTok{ pd}
\CommentTok{\# beetles = pd.read\_csv("resources/data/beetles.csv")}
\CommentTok{\# beetles}
\end{Highlighting}
\end{Shaded}

\begin{tcolorbox}[enhanced jigsaw, coltitle=black, leftrule=.75mm, bottomtitle=1mm, colframe=quarto-callout-tip-color-frame, opacitybacktitle=0.6, colback=white, toprule=.15mm, colbacktitle=quarto-callout-tip-color!10!white, toptitle=1mm, titlerule=0mm, title=\textcolor{quarto-callout-tip-color}{\faLightbulb}\hspace{0.5em}{Conventional tip callout}, arc=.35mm, opacityback=0, rightrule=.15mm, bottomrule=.15mm, breakable, left=2mm]

Some content

\end{tcolorbox}

\phantomsection\label{bestlemma}
\begin{fbxSimple}{Lemma}{Lemma 1: }{Super}
\phantomsection\label{bestlemma}
This is a fantastic Lemma.

\end{fbxSimple}

There are no Lemma's better than Lemma \hyperref[bestlemma]{1}. Note
that referencing of these coloured boxes works like LaTeX's
\texttt{\textbackslash{}ref\{\}} command, thus it just renders the
number, and you have to write the object name in front. This is in
contrast to the \texttt{@lem-name} Quarto notation.

The other way to embed a video is inside a fenced div and treat it like
a figure, as we saw in other chapters.

Our video in this chapter was embedded inside a custom numbered block
and it was called \hyperref[vidFirst]{Video \hyperref[vidFirst]{1}}.

\section{Demonstration of Custom Numbered
Blocks}\label{demonstration-of-custom-numbered-blocks}

Below are examples of the custom-numbered blocks.\\
Each block is foldable (open or closed by default) and has its own
colour scheme and group style as defined in your YAML configuration.

Notice Theorem, Corollary and Lemma has been defined to share counter.
Unlike Video and Definition have their own types (\emph{deflike} and
\emph{vidlike}).

\phantomsection\label{rthmpythag}
\begin{fbxSimple}{Theorem}{Theorem 2: }{Pythagoras’ Theorem}
\phantomsection\label{rthmpythag}
In a right-angled triangle\ldots{}

\end{fbxSimple}

\phantomsection\label{Rcor-pythag}
\begin{fbxSimple}{Corollary}{Corollary 3: }{Corollary to Pythagoras’ Theorem}
\phantomsection\label{Rcor-pythag}
In a right-angled triangle, if the two shorted sides are equal, so that
\(a = b\), then the hypotenuse satisfies \(c=a\sqrt{2}\).

\end{fbxSimple}

\phantomsection\label{Rlem-euclid}
\begin{fbxSimple}{Lemma}{Lemma 4: }{Euclid}
\phantomsection\label{Rlem-euclid}
If two lines are perpendicular to the same line, then they are parallel
to each other.

\end{fbxSimple}

\phantomsection\label{Rvid-yt}
\begin{fbxSimple}{Video}{Video 2: }{Video Demonstration}
\phantomsection\label{Rvid-yt}
Here's a Youtube video:

\url{https://youtu.be/d2nG7Y17sQw}

\end{fbxSimple}

\phantomsection\label{Rdef-orthog}
\begin{fbxSimple}{Definition}{Definition 1: }{Definition of Orthogonality}
\phantomsection\label{Rdef-orthog}
Two vectors \(\mathbf{u}\) and \(\mathbf{v}\) are said to be
\emph{orthogonal} if their dot product is zero: \[
\mathbf{u} \cdot \mathbf{v} = 0.
\]

\end{fbxSimple}

\begin{center}\rule{0.5\linewidth}{0.5pt}\end{center}

\section{Referencing the Blocks}\label{referencing-the-blocks}

You can reference the blocks defined above using their labels:

\begin{itemize}
\tightlist
\item
  The main result is stated in Theorem \hyperref[rthmpythag]{2}.
\item
  Checkout Corollary \hyperref[Rcor-pythag]{3}.\\
\item
  The proof uses Lemma \hyperref[Rlem-euclid]{4}.\\
\item
  Then there's Definition \hyperref[Rdef-orthog]{1}.\\
\item
  Finally, see Video \hyperref[Rvid-yt]{2}.
\end{itemize}

If we want we can wrap the \texttt{\textbackslash{}ref\{\}} inside a
hyperlink to get \hyperref[rthmpythag]{Theorem
\hyperref[rthmpythag]{2}}, \hyperref[Rcor-pythag]{Corollary
\hyperref[Rcor-pythag]{3}}, etc\ldots{}

\section{Customization}\label{customization}

Each box has to define four colours. In the file mentioned above the
light-mode border and background are set. For those using dark-mode you
also need to define a dark mode pair. See
\href{../themes/dark-styles-boxes.scss}{dark styles boxes stylesheet}.

\bookmarksetup{startatroot}

\chapter{Chapter 5}\label{chapter-5}

\section{\texorpdfstring{R package:
\texttt{webexercises}}{R package: webexercises}}\label{r-package-webexercises}

Some people are using external R packages like \texttt{webexercises},
they will just work out of the box. You load it like this, I've added an
extra function to attach accessibility aria labels to the entry boxes
later.

\begin{Shaded}
\begin{Highlighting}[]
\InformationTok{\textasciigrave{}\textasciigrave{}\textasciigrave{}\{r\}}
\FunctionTok{library}\NormalTok{(webexercises)}

\CommentTok{\# To allow injection of aria{-}labels}
\NormalTok{in\_labelled\_fitb }\OtherTok{\textless{}{-}} \ControlFlowTok{function}\NormalTok{(answer, label\_text, ...) \{}
\NormalTok{  html }\OtherTok{\textless{}{-}}\NormalTok{ webexercises}\SpecialCharTok{::}\FunctionTok{fitb}\NormalTok{(answer, ...)}
  \CommentTok{\# Insert aria{-}label into the \textless{}input ...\textgreater{} tag}
\NormalTok{  html }\OtherTok{\textless{}{-}} \FunctionTok{sub}\NormalTok{(}
    \StringTok{"\textless{}input "}\NormalTok{,}
    \FunctionTok{paste0}\NormalTok{(}\StringTok{"\textless{}input aria{-}label=}\SpecialCharTok{\textbackslash{}"}\StringTok{"}\NormalTok{, label\_text, }\StringTok{"}\SpecialCharTok{\textbackslash{}"}\StringTok{ "}\NormalTok{),}
\NormalTok{    html}
\NormalTok{  )}
\NormalTok{  html}
\NormalTok{\}}
\InformationTok{\textasciigrave{}\textasciigrave{}\textasciigrave{}}
\end{Highlighting}
\end{Shaded}

The \texttt{webexercises} package also comes with a css and js file, so
you typically need to add the following to the \texttt{\_metadata.yml}
in the folder for the file being rendered (or in the global
\texttt{\_quarto.yml}) if you want it literally everywhere.

\begin{verbatim}
format:
  html:
    css:
      - include/webex.css
    include-after-body: include/webex.js
\end{verbatim}

Here's an example someone may recognise.

\phantomsection\label{Example-1}
\begin{fbxSimple}{Example}{Example 1}{}
\phantomsection\label{Example-1}
Consider a function taking arguments \texttt{a}, \texttt{b} and
\texttt{c} and returning \((a+2\times b)^c\), defined as set out below.

\begin{Shaded}
\begin{Highlighting}[]
\NormalTok{func }\OtherTok{\textless{}{-}}  \ControlFlowTok{function}\NormalTok{(a, }\AttributeTok{b =} \DecValTok{0}\NormalTok{, }\AttributeTok{c =} \DecValTok{1}\NormalTok{)\{}
\NormalTok{    (a }\SpecialCharTok{+} \DecValTok{2} \SpecialCharTok{*}\NormalTok{ b)}\SpecialCharTok{\^{}}\NormalTok{c}
\NormalTok{\}}
\end{Highlighting}
\end{Shaded}

Rewrite the calls below into named form (i.e.~using
\texttt{parameter\ =\ value}) and give the output of \texttt{func}
without running it in R.

\begin{Shaded}
\begin{Highlighting}[]
\FunctionTok{func}\NormalTok{(}\DecValTok{1}\NormalTok{, }\DecValTok{2}\NormalTok{, }\DecValTok{3}\NormalTok{)}
\end{Highlighting}
\end{Shaded}

This is the equivalent to \_\_\_\_\_\_\_\_\_\_\_\_\_\_\_\_\_ and it
returns the value \_\_\_.

\end{fbxSimple}

\section{Numbas embedding}\label{numbas-embedding}

Another opportunity with HTML is to embed via an iframe to some external
content.

I have embedded some functions insider \texttt{\_common.R} which setup
how to do this via an R function. You need to first create your Numbas
question, then preview it, and click \emph{Download}. This provides a
zip file, which you unzip and place the folder into
\texttt{/resources/numbas/\textless{}your-chosen-name\textgreater{}}.
Then in the code below you just point to the \texttt{index.html} for the
question. Questions are of varying heights and widths too, so you can
customize it in when you create the iframe.

\subsection{Plain version (no web
version)}\label{plain-version-no-web-version}

\begin{Shaded}
\begin{Highlighting}[]
\InformationTok{\textasciigrave{}\textasciigrave{}\textasciigrave{}\{r\}}
\FunctionTok{numbaselement}\NormalTok{(}\StringTok{"/resources/numbas/Q13/index.html"}\NormalTok{, }\AttributeTok{iwidth=}\StringTok{"90\%"}\NormalTok{, }\AttributeTok{iheight=}\StringTok{"900px"}\NormalTok{)}
\InformationTok{\textasciigrave{}\textasciigrave{}\textasciigrave{}}
\end{Highlighting}
\end{Shaded}

A Numbas question appears here in the HTML version.

\subsection{Version with a hyperlink in the
PDF}\label{version-with-a-hyperlink-in-the-pdf}

If you want the PDF version to contain an actual hyperlink then you need
to get a link to a website version of your question, like this example
below. You could try and always use this link rather than the local
download below, but the you are dependent on the external website.

\begin{Shaded}
\begin{Highlighting}[]
\InformationTok{\textasciigrave{}\textasciigrave{}\textasciigrave{}\{r\}}
\FunctionTok{numbaselement}\NormalTok{(}
  \StringTok{"/resources/numbas/Q13/index.html"}\NormalTok{,}
  \AttributeTok{iwidth =} \StringTok{"90\%"}\NormalTok{,}
  \AttributeTok{iheight =} \StringTok{"900px"}\NormalTok{,}
  \AttributeTok{url =} \StringTok{"https://numbas.mathcentre.ac.uk/question/119566/matrix{-}lin\textbackslash{}}
\StringTok{  ear{-}combination/embed/?token=e221efa1{-}3a50{-}429c{-}bb4a{-}8c63d2610114"}
\NormalTok{)}
\InformationTok{\textasciigrave{}\textasciigrave{}\textasciigrave{}}
\end{Highlighting}
\end{Shaded}

\href{https://numbas.mathcentre.ac.uk/question/119566/matrix-lin
  ear-combination/embed/?token=e221efa1-3a50-429c-bb4a-8c63d2610114}{Link to Numbas question}

\begin{tcolorbox}[enhanced jigsaw, coltitle=black, leftrule=.75mm, bottomtitle=1mm, colframe=quarto-callout-note-color-frame, opacitybacktitle=0.6, colback=white, toprule=.15mm, colbacktitle=quarto-callout-note-color!10!white, toptitle=1mm, titlerule=0mm, title=\textcolor{quarto-callout-note-color}{\faInfo}\hspace{0.5em}{Note}, arc=.35mm, opacityback=0, rightrule=.15mm, bottomrule=.15mm, breakable, left=2mm]

I inserted the \texttt{\textbackslash{}} symbol in the long url address,
to render nicely for you, but users don't normally see it.

When presenting you will set \texttt{echo:\ false} on the block anyway,
so the user only sees the question.

\end{tcolorbox}

\subsection{Hiding a numbas task}\label{hiding-a-numbas-task}

An alternative to having a large space-using question, is to hide inside
a collapsed box like the following.

\phantomsection\label{Task-1}
\begin{fbxSimple}{Task}{Task 1}{}
\phantomsection\label{Task-1}

A Numbas question appears here in the HTML version.

\end{fbxSimple}

\bookmarksetup{startatroot}

\chapter{Chapter 6}\label{chapter-6}

\section{Invisible text}\label{invisible-text}

I have a Quarto extension
\href{https://github.com/david-hodge/invis}{Invisible text extension} to
make hiding text easy. Mostly for us in gapped notes, or perhaps
handouts.

It implemented by simply setting text to fully transparent (and not
selectable). Screenreaders \textbf{will} be able to see it, and they
will also get informative notifications about the visibly hidden text.
In LaTeX/PDF it is rendered as white text.

Just install it as any other extension via

\begin{verbatim}
quarto add david-hodge/invis
\end{verbatim}

and then you can use it via usage of the following shortcodes.

\subsection{Simple hiding}\label{simple-hiding}

\begin{description}
\tightlist
\item[\texttt{\{\{\textless{}\ invis\ start\ \textgreater{}\}\}}]
This starts invisibility
\item[\texttt{\{\{\textless{}\ invis\ end\ \textgreater{}\}\}}]
This ends invisibility
\end{description}

By default these are designed to hide a single block/element on a page.
It can hide a single equation, a single word, a single line or text,
etc.. \emph{However} it cannot hide multiple paragraphs or more complex
elements.

\subsection{Hiding larger elements}\label{hiding-larger-elements}

To hide more complex elements, use the \texttt{block} extra paramater.

\begin{description}
\tightlist
\item[\texttt{\{\{\textless{}\ invis\ start\ block\ \textgreater{}\}\}}]
This starts invisibility for a large element
\item[\texttt{\{\{\textless{}\ invis\ end\ block\ \textgreater{}\}\}}]
This ends invisibility for a large element
\end{description}

\subsection{Examples}\label{examples}

This first sentence is fully visible.

\begingroup\color{white} This second sentence is fully visible.
\endgroup

This third sentence is fully visible.

\begin{tcolorbox}[enhanced jigsaw, coltitle=black, leftrule=.75mm, bottomtitle=1mm, colframe=quarto-callout-warning-color-frame, opacitybacktitle=0.6, colback=white, toprule=.15mm, colbacktitle=quarto-callout-warning-color!10!white, toptitle=1mm, titlerule=0mm, title=\textcolor{quarto-callout-warning-color}{\faExclamationTriangle}\hspace{0.5em}{Warning}, arc=.35mm, opacityback=0, rightrule=.15mm, bottomrule=.15mm, breakable, left=2mm]

I don't imagine you would wish, but should you wish to hide some
fundamental Quarto object like a code block and the resulting output you
are recommended to use the built-in features like setting
\texttt{echo:\ false} on the code block.

You can use this \texttt{invis} extension to hide large blocks of
content in the HTML, but for such elements they will likely appear in
the PDF, since all we are doing is setting standard text as white, and
code and plots are not standard text.

\begingroup\color{white}

\begin{Shaded}
\begin{Highlighting}[]
\FunctionTok{plot}\NormalTok{(cars)}
\end{Highlighting}
\end{Shaded}

\begin{figure}[H]

\centering{

\pandocbounded{\includegraphics[keepaspectratio]{Ch6_files/figure-pdf/fig-cars-plot-1.pdf}}

}

\caption{\label{fig-cars-plot}This is the caption}

\end{figure}%

\endgroup

The code and output plot above are fully hidden in the HTML. Likely
fully visible in the PDF.

\end{tcolorbox}




\end{document}
